% ============================================================
\subsection*{The Chain Rule Holds for Every Permutation of the Coordinates}

\subsubsection*{What the statement means}

The chain rule is the definition of conditional probability applied repeatedly:
\[
p(a, b) = p(a)\,p(b \mid a) = p(b)\,p(a \mid b).
\]
Nothing in this identity says \emph{which} variable goes first. So with $D$ variables you can choose any ordering $\pi$ (a permutation of $1,\dots,D$), reveal the variables in that order, and the product of conditionals still equals the same joint:
\[
p(\mathbf{x}) = \prod_{i=1}^{D} p\big(x_{\pi(i)} \mid x_{\pi(1)}, \dots, x_{\pi(i-1)}\big).
\]
Here $\pi(i)$ means ``the coordinate visited at step $i$.''

\subsubsection*{Text example}

Take the 3-token sentence ``the cat sat'', with $x_1=$ the, $x_2=$ cat, $x_3=$ sat.

\textbf{Left-to-right}, $\pi = (1,2,3)$, which is how GPT works:
\[
p(\text{the})\cdot p(\text{cat}\mid\text{the})\cdot p(\text{sat}\mid\text{the, cat})
\]

\textbf{Right-to-left}, $\pi = (3,2,1)$:
\[
p(\text{sat})\cdot p(\text{cat}\mid\text{sat})\cdot p(\text{the}\mid\text{cat, sat})
\]

\textbf{Middle-first}, $\pi = (2,1,3)$:
\[
p(\text{cat})\cdot p(\text{the}\mid\text{cat})\cdot p(\text{sat}\mid\text{the, cat})
\]

All three products give exactly the same number, $p(\text{the, cat, sat})$. The individual factors are different questions, though:
\begin{itemize}
  \item $p(\text{cat} \mid \text{the})$ asks which noun follows ``the''.
  \item $p(\text{cat} \mid \text{sat})$ asks what kind of subject sits.
\end{itemize}

\subsubsection*{Image example}

For a $32\times32$ image, raster order predicts the bottom-right pixel last, using the whole image above it. If you reverse the order, the bottom-right pixel comes \emph{first} and has no context, so its conditional is just the marginal distribution of corner pixel colors. A column-major order or a space-filling curve gives yet another set of conditionals. The image's joint probability is identical under all of them.

\subsubsection*{Why it matters: the order is part of the model}

\begin{enumerate}
  \item \textbf{In theory, the joint is the same.} If the network learned every conditional perfectly, every order would give the same $p(\mathbf{x})$.
  \item \textbf{In practice, the learning problem changes.} The network $f_\theta$ has to approximate a different set of functions for each order. Some orders are easier: text left-to-right matches how grammar unfolds, and audio forward in time matches causality. For images, no order is naturally correct; raster order is a convention.
  \item \textbf{Learned models are not order-invariant.} Two trained models with different orders usually assign \emph{different} probabilities to the same image, because each one's approximation errors fall in different places. Whichever order makes the conditionals easier to learn tends to get better bpd.
  \item \textbf{Sampling follows the order.} Ancestral sampling must generate in order $\pi$. That decides which parts of the output are produced with little context (and are most uncertain), and where exposure-bias errors start to compound.
\end{enumerate}

In one line: the chain rule guarantees that every order is \emph{exactly correct} as a factorization, but you still have to pick one. That choice determines which conditionals the network must learn and the order in which samples are generated.

% ============================================================
\subsection*{Each Order Defines a Different Set of Functions for $f_\theta$}

\subsubsection*{Each conditional is a function}

In an AR model, every factor $p(x_i \mid x_{<i})$ is a \textbf{function}. Its input is the context (the variables already revealed) and its output is a probability distribution over the next variable. The network $f_\theta$, with one set of shared weights $\theta$, has to compute all $D$ of these functions.

The order $\pi$ fixes two things for each step:
\begin{itemize}
  \item \textbf{which variables are the input} (what the network can see), and
  \item \textbf{which variable is the target} (what it must predict).
\end{itemize}
Change the order and both change, so the network is being asked to compute different input$\to$output mappings.

\subsubsection*{Text example}

The functions the network must implement for ``the cat sat'' under each order:

\textbf{Left-to-right}, $\pi=(1,2,3)$:

\begin{tabular}{lll}
\toprule
Step & Input & Output: distribution over \ldots \\
\midrule
1 & nothing & first word of a sentence \\
2 & ``the'' & word after ``the'' \\
3 & ``the cat'' & word after ``the cat'' \\
\bottomrule
\end{tabular}

\textbf{Right-to-left}, $\pi=(3,2,1)$:

\begin{tabular}{lll}
\toprule
Step & Input & Output: distribution over \ldots \\
\midrule
1 & nothing & last word of a sentence \\
2 & ``sat'' & word before ``sat'' \\
3 & ``cat sat'' & word before ``cat sat'' \\
\bottomrule
\end{tabular}

These are different jobs:
\begin{itemize}
  \item Step 1 in left-to-right learns ``sentences often start with \emph{the, I, a}''. In right-to-left, step 1 must learn ``sentences often end with \emph{it, here, sat}''.
  \item Step 2 in left-to-right learns that a noun usually follows a determiner. In right-to-left, it learns which subjects tend to precede a given verb.
\end{itemize}
The data is the same and so is the joint $p(\text{the, cat, sat})$, but the training targets and their inputs are completely different. The weights $\theta$ end up encoding different knowledge.

\subsubsection*{Image example}

Predicting pixel $(16,16)$ in a $32\times32$ face image:
\begin{itemize}
  \item \textbf{Raster order:} the input is the top half plus the left part of row 16. The function maps ``forehead, eyes, part of the nose'' to a distribution over nose-pixel colors. That is fairly easy, because the face structure is already established.
  \item \textbf{Reverse raster:} the input is the bottom half (mouth, chin). The function now maps ``mouth and chin'' to nose color. That is a different function.
  \item \textbf{Center-first order:} this pixel might come first, with no input at all. The function is then just ``what color is the center of any image'', which is very uncertain.
\end{itemize}

\subsubsection*{Why this matters}

\begin{enumerate}
  \item \textbf{Difficulty differs.} Some functions are smooth and easy to learn, others are hard. Left-to-right text suits the way language is structured. For images, every order puts some pixels in a low-context, hard-to-predict spot.
  \item \textbf{Architecture must match.} The masks in PixelCNN and the causal mask in Transformers encode \emph{which inputs each function may see}. Change the order and the masks must change too.
  \item \textbf{Errors differ.} With finite capacity and data, $f_\theta$ approximates its set of functions imperfectly. Different sets lead to different errors, so two models trained with different orders end up with different learned joints $p_\theta(\mathbf{x})$, even though the true joint factors exactly under both.
\end{enumerate}

\textbf{In short:} the order decides what question the network answers at each step (``given \emph{this}, predict \emph{that}''). A different order is a different question set, so the network is a different function approximator, even when it targets the same joint distribution.

% ============================================================
\subsection*{Does a Generative Model Learn Multiple Orders? (Fixed-Order vs.\ Any-Order Models)}

A common misreading is that a generative model learns many different orders of an image so that it can handle any input at inference. This is not how standard AR models work.

\subsubsection*{A standard AR model learns one fixed order}

PixelCNN, PixelRNN and GPT each choose \textbf{one order before training}: raster order for images, left-to-right for text. They learn only the conditionals for that order and never see any other.
\begin{itemize}
  \item \textbf{Training:} for every training image, the network learns ``given pixels $1,\dots,i-1$ in raster order, predict pixel $i$''. It learns nothing about reverse order, column order, or anything else.
  \item \textbf{Inference:} it generates in that same raster order, one pixel after another.
\end{itemize}
The earlier point about different function sets was a comparison \emph{between models}: a separate model trained with a different order would learn a different set of functions. One model does not learn several orders.

\subsubsection*{Generalizing to new inputs is a separate idea}

A standard AR model generalizes because it has seen \textbf{many images in the same order}, not because it has seen many orders. When it meets a new context at inference, such as a partly generated face it has never seen, it relies on patterns learned across the training set, for example ``skin-colored pixels above usually mean more skin below''.

There is one limit. It can only condition on information that comes \textbf{earlier in its order}. Inpainting works naturally because known pixels above and to the left act as context. Pixels below or to the right of the hole are ignored by a raster model, even when they are known, because they come later in its order.

\subsubsection*{Models that do learn many orders}

This idea exists as a separate family, \textbf{order-agnostic} or \textbf{any-order} AR models:
\begin{itemize}
  \item \textbf{Order-agnostic MADE} and \textbf{XLNet} sample a random permutation $\pi$ for each training example, so one network learns conditionals for many orders.
  \item The benefit is that at inference you can choose the order, for example generating the unknown pixels conditioned on all known pixels wherever they are.
  \item The cost is a much harder learning problem (many function sets in one network), which usually means worse likelihood than a fixed-order model of the same size.
\end{itemize}

\begin{tabular}{lll}
\toprule
 & Standard AR (PixelCNN, GPT) & Any-order AR (XLNet, MADE) \\
\midrule
Orders learned & 1, fixed & many, random permutations \\
Generalizes by & many images in that order & many images $\times$ many orders \\
Inference order & must follow the trained order & can choose \\
\bottomrule
\end{tabular}

% ============================================================
\subsection*{Directed Graph of an Autoregressive Model}

\subsubsection*{What the graph shows}

For four variables, the graph has these arrows:
\[
x_1 \to x_2,\quad x_1 \to x_3,\quad x_1 \to x_4,\quad x_2 \to x_3,\quad x_2 \to x_4,\quad x_3 \to x_4.
\]
Each \textbf{node} is one variable: a pixel, a token, or an audio sample. Each \textbf{arrow} $x_j \to x_i$ means ``$x_i$ may depend on $x_j$''. The variables that point into a node are its \textbf{parents}, and the parents are exactly the context in the chain rule:

\begin{tabular}{lll}
\toprule
Node & Parents & Factor \\
\midrule
$x_1$ & none & $p(x_1)$ \\
$x_2$ & $x_1$ & $p(x_2 \mid x_1)$ \\
$x_3$ & $x_1, x_2$ & $p(x_3 \mid x_1, x_2)$ \\
$x_4$ & $x_1, x_2, x_3$ & $p(x_4 \mid x_1, x_2, x_3)$ \\
\bottomrule
\end{tabular}

Every earlier node points to every later node. This is a \textbf{fully connected} forward graph, which is why an AR model makes no independence assumptions: the chain rule itself is the graph.

\textbf{Audio example:} with $x_1,\dots,x_4$ as four consecutive waveform samples, the sample at time 4 may depend on all three previous samples, but sample 2 cannot depend on sample 3.

\subsubsection*{Why ``no cycle'' matters}

All arrows point forward, so the graph is a \textbf{DAG} (directed acyclic graph). That guarantees a valid generation order: start at the node with no parents and move forward. This is ancestral sampling:
\begin{enumerate}
  \item Sample the first token, ``the'', from $p(x_1)$.
  \item Sample ``cat'' from $p(x_2 \mid \text{the})$.
  \item Continue to the end.
\end{enumerate}
If there were a cycle, say $x_2$ depends on $x_3$ and $x_3$ depends on $x_2$, neither could be generated first, because each would wait for the other. Sampling would be undefined.

\subsubsection*{One network with shared weights approximates all conditionals}

The graph has $D$ different conditionals, each with a different number of parents. You don't train $D$ separate networks. One network $f_\theta$ computes all of them, and masking controls which inputs each output position can see. A Transformer predicting token 4 and token 2 uses the same weights; only the visible context differs.

\subsubsection*{The causal constraint}

Predicting $x_i$ must not depend on $x_i, x_{i+1}, \dots, x_D$. Graphically, \textbf{no arrow may point backward or into a node from itself}. The product is only a valid probability if every factor uses earlier variables alone.

\textbf{What goes wrong if it leaks:} suppose a text model predicting token $x_i$ can see $x_i$ itself. It learns to copy: $p(x_i \mid \ldots, x_i) = 1$ for every token. Then for \emph{any} sentence,
\[
\prod_i p(\cdot) = 1 \times 1 \times \cdots = 1.
\]
Every possible sentence receives probability 1, and the total over all sentences becomes far greater than 1. It is no longer a distribution, so the likelihood and bpd are meaningless. The model looks perfect during training but generates garbage, because at sampling time the ``answer'' it learned to copy is not available yet.

\subsubsection*{Where masking comes in}

A CNN or Transformer naturally looks at all positions, which would add backward arrows. Masking deletes those arrows:
\begin{itemize}
  \item \textbf{PixelCNN:} Mask A/B zero out the filter weights that touch the current pixel or future pixels.
  \item \textbf{GPT:} the causal attention mask sets future positions to $-\infty$.
\end{itemize}
Both enforce exactly this DAG inside the network.